\documentclass[11pt]{article}
\usepackage[margin=1in]{geometry}
\usepackage{amsmath,amssymb,amsthm}
\usepackage{booktabs}
\usepackage{graphicx}
\usepackage{hyperref}
\usepackage{xcolor}
\usepackage{authblk}
\usepackage{natbib}

\title{Learnable Path-Complex Features and Unified Path--Cellular Message
Passing for Link Prediction}

\author{Farhan Sadeek}
\affil{Dartmouth College}
\date{April 2026}

\newcommand{\pcn}{\textsc{PCN}}
\newcommand{\cwn}{\textsc{CWN}}
\newcommand{\upcn}{\textsc{UPCN}}

\begin{document}
\maketitle

\begin{abstract}
  Path Complex Networks (\pcn{}) generalize simplicial and cellular message
  passing by lifting a graph to its path complex and exchanging messages
  between $k$-paths through boundary, coboundary, and upper/lower-adjacency
  relations. In the original formulation, the feature of a $k$-cell
  $(v_0,\ldots,v_k)$ is a deterministic scatter-sum (or scatter-mean) of
  its constituent node features, discarding both the \emph{order} along
  the path and any learnable capacity at the lifting step. Moreover, \pcn{}
  commits to a single topological lift---paths, cliques, or rings---so the
  model cannot simultaneously exploit the complementary inductive biases
  offered by different cell types. We address both limitations. First, we
  introduce a position-aware learnable feature initialization in which
  each position along a $k$-cell receives its own linear (or MLP) map,
  yielding order-sensitive path features at no asymptotic cost. Second,
  we construct a unified path\,+\,cellular complex in which path-cells and
  ring-cells coexist at the same dimension and communicate through shared
  edge boundaries, so that a triangle $ABC$ is simultaneously a path
  $A{\to}B{\to}C$ and a 3-cycle without duplicating adjacency bookkeeping.
  We evaluate both contributions on Cora link prediction. Our experiments
  ablate (i)~deterministic vs.\ learnable initialization and (ii)~path-only
  vs.\ unified complexes, and we report test AUC across 10 seeds for every
  configuration.
\end{abstract}

%-----------------------------------------------------------------------
\section{Introduction}

Graph neural networks that rely on the 1-Weisfeiler--Lehman (1-WL)
update rule are fundamentally limited in their ability to distinguish
non-isomorphic graphs \citep{xu2019powerful}. Higher-order topological
GNNs escape this limit by lifting the graph to a richer combinatorial
object---a simplicial complex \citep{bodnar2021mpsn}, a regular cell
complex \citep{bodnar2021cwn}, or a path complex
\citep{truong2024pcn}---and running message passing over its cells.
Each lift carries a different inductive bias: simplicial complexes
privilege cliques, cell complexes privilege induced cycles, and path
complexes privilege simple paths and thus make no structural
assumption at all.

\pcn{} \citep{truong2024pcn} is attractive precisely because paths
exist in every graph, but the published implementation leaves two
degrees of freedom unused. (1)~Cell features at dimension $k\ge1$ are
initialized by a deterministic scatter-sum of node features. Because
summation is symmetric, the feature of the path $A{\to}B{\to}C$ is
identical to the feature of $C{\to}B{\to}A$, which is counterintuitive
for an object whose defining property is an ordering of its vertices.
(2)~The lift is monolithic: one picks paths \emph{or} rings \emph{or}
cliques. In practice, however, the same substructure---say a
triangle---is a clique, a 3-cycle, and a path of length two
simultaneously, and a model with access to all three views should, in
principle, be strictly more expressive than any one of them.

We propose two small but targeted modifications to \pcn{}:

\begin{enumerate}
  \item \textbf{Learnable, position-aware cell initialization.}
    For a $k$-cell $(v_0,\ldots,v_k)$ we learn $k{+}1$ linear
    maps $W_0,\ldots,W_k$, one per position, and define
    $\mathbf{h}^{(0)}_{\sigma} = \sum_{i=0}^{k} W_i\mathbf{x}_{v_i}$.
    The node features now enter the path-complex features through
    a learned, ordered transform rather than a symmetric sum.
  \item \textbf{Unified path\,+\,cellular complex.}
    We build a single complex in which path-cells and ring-cells
    coexist at dimension two, share the same edge id-map, and
    communicate through their shared edges during message passing.
    A triangle therefore sends messages \emph{as} a ring-cell and
    simultaneously appears as a path-cell; the two views are
    fused downstream through the dim-1 edges they both bound.
\end{enumerate}

Both modifications are drop-in: the message-passing layer is unchanged,
and the per-cell feature dimension is unchanged. We evaluate on Cora
link prediction in the transductive setting of
\citet{kipf2016variational}, reporting mean $\pm$ std test AUC over
10 seeds for four configurations: (path,\,det), (path,\,learn),
(unified,\,det), (unified,\,learn).

%-----------------------------------------------------------------------
\section{Background}

\paragraph{Path complexes.}
Given a graph $G=(V,E)$, a \emph{simple path} of length $k$ is an
ordered tuple $(v_0,\ldots,v_k)$ of distinct vertices with
$\{v_i,v_{i+1}\}\in E$ for all $i$. The set $P_k(G)$ of all such
paths, together with boundary maps $\partial$ defined by deletion of
the leftmost or rightmost vertex, forms a \emph{path complex}
\citep{grigoryan2012homologies,truong2024pcn}. Message passing in
\pcn{} operates over $P_0(G)\cup P_1(G)\cup\cdots\cup P_K(G)$ via
boundary, coboundary, and upper/lower adjacencies between paths, in
direct analogy with \cwn{}.

\paragraph{Deterministic feature initialization.}
In the public \pcn{} implementation, the feature of a $k$-cell at the
input layer is
\begin{equation}
  \mathbf{h}^{(0)}_{\sigma} = \bigoplus_{v\in\sigma} \mathbf{x}_v,
  \qquad \sigma=(v_0,\ldots,v_k),
  \label{eq:det-init}
\end{equation}
where $\oplus\in\{\text{sum},\text{mean}\}$. We refer to this as the
\emph{deterministic init}. Equation~\eqref{eq:det-init} is symmetric
under any permutation of $(v_0,\ldots,v_k)$, so path orientation is
lost and cannot be recovered by subsequent layers that operate only on
$\mathbf{h}^{(0)}_\sigma$.

%-----------------------------------------------------------------------
\section{Method}

\subsection{Learnable position-aware initialization}

We replace \eqref{eq:det-init} with
\begin{equation}
  \mathbf{h}^{(0)}_{\sigma}
  = \sum_{i=0}^{k} W_i^{(k)}\,\mathbf{x}_{v_i},
  \qquad \sigma=(v_0,\ldots,v_k),
  \label{eq:learn-init}
\end{equation}
where $W_i^{(k)}\in\mathbb{R}^{d\times d_\text{in}}$ is a learnable
linear map specific to the pair $(k,i)$ of dimension and position.
Asymptotically the cost is identical to \eqref{eq:det-init}: each node
contributes exactly once to the feature of each cell it bounds, and
the only change is that the contribution is multiplied by a
position-specific matrix before being accumulated. When $W_i^{(k)}=W$
for all $(k,i)$, \eqref{eq:learn-init} reduces to a linear projection
of the scatter-sum, so the deterministic init is strictly less
expressive.

In the implementation, the path-complex builder already records a
\texttt{cell\_node\_index} tensor of shape $[2,\sum_\sigma (k_\sigma+1)]$
and a \texttt{position\_index} tensor; the learnable init iterates over
positions $0,\ldots,k$ and performs a masked linear projection followed
by an \texttt{index\_add\_} to accumulate into the cell feature buffer
(\texttt{lib/layers/learnable\_features.py}).

\subsection{Unified path--cellular complex}

Let $P_k(G)$ denote the set of simple $k$-paths and $R(G)$ the set of
induced cycles up to length $L$. We build a \emph{unified complex}
$\mathcal{C}(G)$ with cells
\[
  \mathcal{C}_0 = V,\quad
  \mathcal{C}_1 = P_1(G),\quad
  \mathcal{C}_2 = P_2(G)\ \cup\ R(G),\quad
  \mathcal{C}_k = P_k(G)\ \text{for } k\ge3.
\]
The dim-0 and dim-1 id-maps are shared across path-cells and ring-cells
(edges are canonicalized to sorted tuples), so a ring-cell and a
path-cell that share an edge point to the same edge id. Boundary
extraction branches on cell type: path-cells use the left/right-deletion
rule, and ring-cells use the standard CW boundary map that returns every
edge of the cycle. Upper and lower adjacencies are then computed by the
existing \pcn{} routine on the combined boundary tables, so messages
flow seamlessly between a ring-cell and a path-cell that happen to share
an edge.

\paragraph{Why this helps link prediction.} Link prediction rewards
models whose node embeddings capture closed-neighborhood structure: a
pair $(u,v)$ is more likely to be linked when $u$ and $v$ participate in
short cycles. Path-cells expose only open neighborhoods; ring-cells
expose closed ones. Since cycles of length $\le L$ in Cora are
overwhelmingly short, adding ring-cells to a path complex injects a
targeted topological signal that paths alone do not carry.

%-----------------------------------------------------------------------
\section{Experiments}

\subsection{Setup}

We evaluate on Cora in the transductive link-prediction setting of
\citet{kipf2016variational}: edges are split $85/5/10$ into
train/val/test with an equal number of sampled negatives per split.
The path complex (or unified complex) is built from training
message-passing edges only. All models use a 2-layer \pcn{} encoder
with hidden dimension 64, dropout 0.3, batch normalization, a dot-product
decoder, and Adam with $\text{lr}=0.01$ and
\texttt{ReduceLROnPlateau} scheduling. Each run is 200 epochs with
early stopping (patience 20 on val AUC). We report mean $\pm$ std over
10 seeds ($0,\ldots,9$).

\subsection{Ablations}

\begin{table}[t]
  \centering
  \small
  \begin{tabular}{llcc}
    \toprule
    Complex & Init & Val AUC & Test AUC \\
    \midrule
    Path      & Deterministic & $0.7344 \pm 0.0274$ & $0.7113 \pm 0.0239$ \\
    Path      & Learnable     & \textcolor{gray}{[to be filled from \texttt{pcn\_cora\_lp\_learnable}]} & \textcolor{gray}{[\ldots]} \\
    Unified   & Deterministic & \textcolor{gray}{[to be filled from \texttt{pcn\_cora\_lp\_unified}]} & \textcolor{gray}{[\ldots]} \\
    Unified   & Learnable     & \textcolor{gray}{[to be filled from \texttt{pcn\_cora\_lp\_unified\_learnable}]} & \textcolor{gray}{[\ldots]} \\
    \bottomrule
  \end{tabular}
  \caption{Cora link-prediction AUC (mean $\pm$ std over 10 seeds).
    Baseline (Path, Deterministic) numbers are from the
    \texttt{pcn\_cora\_link\_prediction} run at
    \texttt{experiments/results/Cora/runs/2026-04-05T13:11:20-\ldots}.
    The remaining three rows are produced by the three scripts in
  \texttt{experiments/scripts/}.}
  \label{tab:ablation}
\end{table}

Table~\ref{tab:ablation} reports the four configurations. The baseline
(path, deterministic) achieves test AUC $0.7113 \pm 0.0239$, consistent
with the fact that \pcn{} was not designed for link prediction and a
pure-path lift provides little closed-neighborhood signal. We compare
against this baseline the three modifications introduced in
Section~3.

\subsection{Discussion}

Neither of our modifications is expected to close the gap to
specialised link-prediction architectures such as SEAL
\citep{zhang2018seal}, NBFNet, or a well-tuned VGAE---those models
exploit properties of the link-prediction problem (subgraph
  enclosure, Bellman--Ford relational paths, bilinear decoders on
structural features) that are largely orthogonal to the topological
lift studied here. What the ablation does isolate is the marginal
effect, \emph{within} the \pcn{} family, of (i)~learning the lift from
data and (ii)~admitting multiple cell types simultaneously.

A natural limitation of the unified complex is duplication: a
3-cycle $ABC$ is now represented both as a ring-cell and as three
path-cells $A{\to}B{\to}C$, $B{\to}C{\to}A$, $C{\to}A{\to}B$. We do
not deduplicate, and we expect this to over-weight triangle
structure; exploring learnable gating between cell types is a
promising direction.

%-----------------------------------------------------------------------
\section{Conclusion}

We described two small extensions to Path Complex Networks---learnable
position-aware cell initialization and a unified path\,+\,cellular
lift---and evaluated them on Cora link prediction. Both are
implemented as drop-in modifications to the public \pcn{} codebase,
require no change to the message-passing layer, and preserve the
feature budget of the original model. Future work includes extending
the unified complex to admit clique-cells (simplicial) alongside
paths and rings, evaluating on OGB link-prediction benchmarks, and
testing whether the position-aware init offers benefits for
graph-level tasks where \pcn{} is already competitive.

\bibliographystyle{plainnat}
\begin{thebibliography}{9}
  \bibitem[Bodnar et al.(2021a)]{bodnar2021mpsn}
  Bodnar, C., Frasca, F., Wang, Y.~G., Otter, N., Mont\'ufar, G., Li\`o, P.,
  and Bronstein, M.
  \newblock Weisfeiler and Lehman go topological: Message passing simplicial
  networks.
  \newblock In \emph{ICML}, 2021.

  \bibitem[Bodnar et al.(2021b)]{bodnar2021cwn}
  Bodnar, C., Frasca, F., Otter, N., Wang, Y.~G., Li\`o, P., Mont\'ufar, G.,
  and Bronstein, M.
  \newblock Weisfeiler and Lehman go cellular: CW networks.
  \newblock In \emph{NeurIPS}, 2021.

  \bibitem[Grigor'yan et al.(2012)]{grigoryan2012homologies}
  Grigor'yan, A., Lin, Y., Muranov, Y., and Yau, S.-T.
  \newblock Homologies of path complexes and digraphs.
  \newblock \emph{arXiv:1207.2834}, 2012.

  \bibitem[Kipf and Welling(2016)]{kipf2016variational}
  Kipf, T.~N.\ and Welling, M.
  \newblock Variational graph auto-encoders.
  \newblock In \emph{NeurIPS Bayesian Deep Learning Workshop}, 2016.

  \bibitem[Truong and Chin(2024)]{truong2024pcn}
  Truong, Q.\ and Chin, P.
  \newblock Weisfeiler and Lehman go paths: Learning topological features via
  path complexes.
  \newblock In \emph{AAAI}, 2024.

  \bibitem[Xu et al.(2019)]{xu2019powerful}
  Xu, K., Hu, W., Leskovec, J., and Jegelka, S.
  \newblock How powerful are graph neural networks?
  \newblock In \emph{ICLR}, 2019.

  \bibitem[Zhang and Chen(2018)]{zhang2018seal}
  Zhang, M.\ and Chen, Y.
  \newblock Link prediction based on graph neural networks.
  \newblock In \emph{NeurIPS}, 2018.
\end{thebibliography}

\end{document}
